\section{Related Work}
\label{sec:related-work}
Various works have explored the reuse of the KV cache, and many of them are concurrent with our work. Uniquely, our RadixAttention first proposes treating the KV cache as a tree-based LRU cache. It is the first solution that supports multi-level sharing, cache-aware scheduling, frontend-runtime co-scheduling, and distributed cases.
vLLM~\cite{kwon2023vllm} and ChunkedAttention~\cite{ye2024chunkattention} explore some simple reuse cases (e.g., system prompt sharing) but do not cover multi-level tree-structured sharing or LRU caching. PromptCache~\cite{gim2023prompt} proposes the modular reuse of the KV cache beyond the prefix but can impact accuracy by up to a 43\% drop. HydraGen~\cite{juravsky2024hydragen}, FlashInfer~\cite{flashinfer}, and ChunkedAttention~\cite{ye2024chunkattention} focus on CUDA kernel optimizations and do not include the concept of an LRU cache. API Serve~\cite{abhyankar2024apiserve} and LLM-SQL~\cite{liu2024optimizing} study KV cache reuse for specific applications such as interleaving with external API calls and relational databases, but they do not have our radix tree or cache-aware scheduling.

Several LLM programming and agent frameworks exist, such as Guidance~\cite{guidance}, LMQL~\cite{beurer2023lmql}, DSPy~\cite{khattab2023dspy}, LangChain~\cite{langchain}, AutoGen~\cite{wu2023autogen}, and LLM Compiler~\cite{kim2023llm}. Guidance and LMQL are most similar to \lang, and we compare them in \autoref{sec:programming_model}. Our innovation lies in novel runtime optimizations for accelerating the proposed programming model. \lang is compatible with other frameworks and can accelerate them (e.g., the DSPy example in our evaluation). Additionally, \lang is compatible with many other common inference optimizations ~\cite{yu2022orca, pope2023efficiently, aminabadi2022deepspeed, kwon2023vllm, flashinfer, dao2022flashattention, lin2023awq, hooper2024kvquant, kang2024gear, liu2024kivi, liu2024scissorhands, ge2023model}.

